\chapter{การวิเคราะห์ข้อมูลราสเตอร์ (Raster-based Spatial Analysis)}

\section{วัตถุประสงค์การเรียนรู้}
หลังจากศึกษาจบบทนี้ ผู้เรียนจะสามารถ:
\begin{enumerate}
    \item อธิบายความหมายและความสำคัญของระบบสารสนเทศภูมิศาสตร์แบบราสเตอร์ได้
    \item ดำเนินการ Map Algebra ขั้นพื้นฐานและขั้นสูงได้
    \item จำแนกประเภทการวิเคราะห์เชิงพื้นที่ (Local, Neighborhood, Zonal, Global) และประยุกต์ใช้ได้
    \item คำนวณและแปลความหมายของ Slope, Aspect, Hillshade และ Viewshed ได้
    \item ประยุกต์ใช้การวิเคราะห์ราสเตอร์ในงานด้านทรัพยากรธรรมชาติและสิ่งแวดล้อมได้
    \item แปลงข้อมูลระหว่างรูปแบบ Vector และ Raster สำหรับการวิเคราะห์ได้
\end{enumerate}

\section{บทนำ}
ระบบสารสนเทศภูมิศาสตร์ (Geographic Information System: GIS) เป็นเทคโนโลยีที่ใช้ในการเก็บรวบรวม จัดเก็บ จัดการ วิเคราะห์ และแสดงผลข้อมูลเชิงพื้นที่ ข้อมูลเชิงพื้นที่สามารถจัดเก็บได้สองรูปแบบหลัก ได้แก่ รูปแบบเวกเตอร์ (Vector) ที่ใช้จุด เส้น และรูปหลายเหลี่ยมแทนข้อมูล และรูปแบบราสเตอร์ (Raster) ที่ใช้กริดของพิกเซลในการแทนค่าความเข้มข้นหรือคุณลักษณะของข้อมูลในพื้นที่

ในบทนี้จะกล่าวถึงการวิเคราะห์ข้อมูลราสเตอร์ (Raster-based Spatial Analysis) ซึ่งเป็นวิธีการวิเคราะห์ข้อมูลเชิงพื้นที่ที่ใช้โครงสร้างข้อมูลแบบกริด การวิเคราะห์ราสเตอร์มีความสำคัญอย่างยิ่งในงานด้านทรัพยากรธรรมชาติและสิ่งแวดล้อม เนื่องจากข้อมูลราสเตอร์สามารถแสดงปรากฏการณ์ที่มีการเปลี่ยนแปลงอย่างต่อเนื่องในพื้นที่ได้อย่างมีประสิทธิภาพ เช่น อุณหภูมิ ความสูง ปริมาณน้ำฝน และการกระจายตัวของพืชพรรณ

การวิเคราะห์ข้อมูลราสเตอร์ประกอบด้วยหัวข้อหลักหลายประการ ได้แก่ Map Algebra ซึ่งเป็นพีชคณิตสำหรับการคำนวณแผนที่ การวิเคราะห์เชิงพื้นที่ที่แบ่งตามขอบเขตการคำนวณ การวิเคราะห์พื้นผิวที่ได้จากข้อมูลระดับสูง และการแปลงข้อมูลระหว่างรูปแบบ Vector และ Raster ความเข้าใจในหลักการเหล่านี้จะทำให้สามารถนำไปประยุกต์ใช้ในงานวิจัยและการปฏิบัติงานจริงได้
\label{sec:raster-intro}

\section{โครงสร้างข้อมูลราสเตอร์}
\index{ราสเตอร์}
\index{Raster}

\subsection{นิยามของข้อมูลราสเตอร์}
ข้อมูลราสเตอร์ (Raster Data) คือ ข้อมูลเชิงพื้นที่ที่จัดเก็บในรูปแบบกริดสองมิติ ประกอบด้วยแถวและคอลัมน์ของพิกเซล (Pixel) หรือเซลล์ (Cell) ที่มีขนาดเท่ากันทุกตัว โดยแต่ละพิกเซลเก็บค่าหนึ่งค่าที่แสดงถึงคุณลักษณะหรือปรากฏการณ์ในพื้นที่ที่พิกเซลนั้นครอบคลุม ขนาดของพิกเซลเรียกว่า ความละเอียดเชิงพื้นที่ (Spatial Resolution) ซึ่งกำหนดรายละเอียดของข้อมูล

การใช้ข้อมูลราสเตอร์มีข้อได้เปรียบหลายประการ ได้แก่ โครงสร้างข้อมูลที่เรียบง่ายและเข้าใจง่าย การคำนวณทางคณิตศาสตร์ที่มีประสิทธิภาพ และความเหมาะสมกับข้อมูลที่มีการเปลี่ยนแปลงอย่างต่อเนื่องในพื้นที่ อย่างไรก็ตาม ข้อมูลราสเตอร์มีข้อจำกัดในเรื่องขนาดไฟล์ที่ใหญ่เมื่อใช้ความละเอียดสูง และความไม่แม่นยำในการแสดงขอบเขตของวัตถุที่มีรูปร่างซับซ้อน

\begin{ex}
พิจารณาข้อมูลราสเตอร์ที่แสดงระดับความสูงของพื้นที่ กริดขนาด $4 \times 4$ พิกเซล ซึ่งแต่ละพิกเซลมีค่าระดับความสูงเป็นเมตร ดังแสดงในตารางที่~\ref{tab:raster-example}

\begin{table}[ht]
\centering
\caption{ตัวอย่างข้อมูลราสเตอร์แสดงระดับความสูง (หน่วย: เมตร)}
\label{tab:raster-example}
\begin{tabular}{|c|c|c|c|c|}
\hline
 & \textbf{คอลัมน์ 1} & \textbf{คอลัมน์ 2} & \textbf{คอลัมน์ 3} & \textbf{คอลัมน์ 4} \\ \hline
\textbf{แถว 1} & 100 & 120 & 135 & 142 \\ \hline
\textbf{แถว 2} & 105 & 118 & 130 & 140 \\ \hline
\textbf{แถว 3} & 95  & 110 & 125 & 138 \\ \hline
\textbf{แถว 4} & 88  & 98  & 115 & 132 \\ \hline
\end{tabular}
\end{table}

ในตารางที่~\ref{tab:raster-example} แต่ละพิกเซลแทนค่าระดับความสูงเฉลี่ยในพื้นที่ที่พิกเซลนั้นครอบคลุม ค่าความสูงเพิ่มขึ้นจากซ้ายไปขวาและจากล่างขึ้นบน ซึ่งบ่งบอกว่าพื้นที่มีความสูงมากทางด้านขวาบน
\end{ex}

\subsection{ประเภทของข้อมูลราสเตอร์}
ข้อมูลราสเตอร์สามารถจำแนกได้หลายประเภทตามลักษณะของค่าที่เก็บในพิกเซล การจำแนกประเภทข้อมูลราสเตอร์มีความสำคัญในการเลือกวิธีการวิเคราะห์ที่เหมาะสม

\begin{enumerate}
    \item ข้อมูลเชิงปริมาณต่อเนื่อง (Continuous Data) คือ ข้อมูลที่ค่าในพิกเซลสามารถเปลี่ยนแปลงได้อย่างต่อเนื่องในช่วงค่าที่กำหนด ตัวอย่างเช่น ระดับความสูง อุณหภูมิ ความชื้น และความเข้มของการสะท้อนแสงจากดาวเทียม ข้อมูลประเภทนี้มักใช้ในการวิเคราะห์พื้นผิวและการประมาณค่าในตำแหน่งที่ไม่มีข้อมูล

    \item ข้อมูลเชิงหมวดหมู่ (Categorical Data) คือ ข้อมูลที่ค่าในพิกเซลแทนประเภทหรือคลาสของวัตถุ เช่น ประเภทการใช้ที่ดิน ชนิดพืชพรรณ และเขตการปกครอง ค่าในข้อมูลประเภทนี้มักเป็นจำนวนเต็มที่แทนรหัสของแต่ละประเภท

    \item ข้อมูลเชิงลำดับ (Ordinal Data) คือ ข้อมูลที่มีการจัดลำดับความสำคัญ โดยค่าที่สูงกว่าหมายถึงระดับที่สูงกว่า แต่ไม่สามารถระบุความแตกต่างเชิงปริมาณได้ ตัวอย่างเช่น ระดับความเสี่ยงของการเกิดไฟป่า (ต่ำ กลาง สูง) หรือระดับความเหมาะสมของที่ดินสำหรับเกษตรกรรม
\end{enumerate}

\section{พีชคณิตแผนที่ (Map Algebra)}
\index{พีชคณิตแผนที่}
\index{Map Algebra}

\subsection{แนะนำ Map Algebra}
Map Algebra เป็นชุดของการดำเนินการทางคณิตศาสตร์และตรรกะที่ใช้ในการวิเคราะห์ข้อมูลราสเตอร์ Map Algebra ถูกพัฒนาขึ้นโดย Dana Tomlin ในทศวรรษ 1980 และกลายเป็นพื้นฐานของการวิเคราะห์ราสเตอร์ในระบบ GIS สมัยใหม่ หลักการพื้นฐานของ Map Algebra คือการปฏิบัติการทางคณิตศาสตร์จะถูกนำไปใช้กับพิกเซลที่ตำแหน่งเดียวกันในกริดต่างๆ ผลลัพธ์จะเป็นกริดใหม่ที่มีค่าของการปฏิบัติการ

Map Algebra มีความสำคัญอย่างยิ่งในการวิเคราะห์เชิงพื้นที่ เนื่องจากเป็นเครื่องมือที่ทรงพลังในการรวมข้อมูลจากหลายแหล่ง คำนวณดัชนีทางวิทยาศาสตร์ และสร้างแบบจำลองเชิงพื้นที่ที่ซับซ้อน การเข้าใจ Map Algebra อย่างถ่องแท้จะเป็นรากฐานสำหรับการเรียนรู้การวิเคราะห์ราสเตอร์ขั้นสูงต่อไป

การดำเนินการ Map Algebra สามารถจำแนกได้เป็น 4 ประเภทหลักตามจำนวนมิติที่ใช้ในการคำนวณ ได้แก่ Local Operations ที่ทำงานกับพิกเซลเดี่ยว Neighborhood Operations ที่ทำงานกับพิกเซลและพิกเซลรอบข้าง Zonal Operations ที่ทำงานกับกลุ่มของพิกเซลที่อยู่ในโซนเดียวกัน และ Global Operations ที่ทำงานกับพิกเซลทั้งหมดในกริด

\subsection{การดำเนินการระดับพิกเซล (Local Operations)}
\index{Local Operations}

การดำเนินการระดับพิกเซล (Local Operations) คือ การดำเนินการที่คำนวณค่าของพิกเซลผลลัพธ์จากค่าของพิกเซลเดี่ยวหนึ่งตัวหรือหลายตัวที่ตำแหน่งเดียวกันในกริดต่างๆ การดำเนินการระดับพิกเซลเป็นพื้นฐานที่สุดของ Map Algebra เนื่องจากทำงานโดยตรงกับค่าของแต่ละพิกเซลโดยไม่ต้องพิจารณาความสัมพันธ์กับพิกเซลรอบข้าง

การดำเนินการระดับพิกเซลสามารถแบ่งได้เป็นการดำเนินการทางคณิตศาสตร์และการดำเนินการทางตรรกะ การดำเนินการทางคณิตศาสตร์ประกอบด้วยการบวก การลบ การคูณ การหาร และฟังก์ชันทางคณิตศาสตร์อื่นๆ การดำเนินการทางตรรกะประกอบด้วย AND OR NOT และ XOR ซึ่งมักใช้ในการรวมข้อมูลเชิงหมวดหมู่

\begin{ex}
การบวกกริดสองชั้นข้อมูล (Grid Addition) พิจารณากริดสองชั้นที่แสดงปริมาณน้ำฝนรายเดือนของเดือนมกราคมและเดือนกุมภาพันธ์ การบวกกริดทั้งสองจะได้กริดผลลัพธ์ที่แสดงปริมาณน้ำฝนรวมของทั้งสองเดือน สำหรับพิกเซลที่ตำแหน่ง $(i, j)$ ค่าผลลัพธ์จะเป็น $R_{ij} = G1_{ij} + G2_{ij}$ โดยที่ $G1_{ij}$ และ $G2_{ij}$ เป็นค่าปริมาณน้ำฝนของเดือนมกราคมและกุมภาพันธ์ตามลำดับ
\end{ex}

\begin{ex}
การคูณกริดด้วยค่าคงที่ (Grid Scalar Multiplication) ในการแปลงหน่วยของข้อมูล อาจต้องคูณกริดด้วยค่าคงที่ ตัวอย่างเช่น หากข้อมูลระดับความสูงเก็บในหน่วยเซนติเมตร แต่ต้องการหน่วยเป็นเมตร จะต้องคูณค่าทุกพิกเซลด้วย $0.01$ สำหรับพิกเซลที่ตำแหน่ง $(i, j)$ ค่าผลลัพธ์จะเป็น $R_{ij} = 0.01 \times G_{ij}$
\end{ex}

\begin{ex}
การดำเนินการทางตรรกะ (Logical Operations) การดำเนินการทางตรรกะมักใช้ในการสร้างแผนที่แสดงพื้นที่ที่ตรงตามเงื่อนไขหลายประการ ตัวอย่างเช่น การหาพื้นที่ที่มีความสูงมากกว่า 1000 เมตร และมีความลาดชันน้อยกว่า 15 องศา สำหรับพิกเซลที่ตำแหน่ง $(i, j)$ ผลลัพธ์จะเป็น $R_{ij} = (H_{ij} > 1000) \land (S_{ij} < 15)$ โดยที่ $H$ คือกริดความสูง และ $S$ คือกริดความลาดชัน
\end{ex}

\begin{ex}
การคำนวณดัชนีพืชพรรณ (Vegetation Index) ดัชนีพืชพรรณปกติ (Normalized Difference Vegetation Index: NDVI) คำนวณจากข้อมูลการสะท้อนแสงในช่วงแสงแดง (Red) และแสงใกล้อินฟราเรด (Near-Infrared: NIR) สูตรคือ

\[
\text{NDVI} = \frac{\text{NIR} - \text{Red}}{\text{NIR} + \text{Red}}
\]

โดยที่ NIR คือค่าการสะท้อนในช่วงใกล้อินฟราเรด และ Red คือค่าการสะท้อนในช่วงแสงแดง ค่า NDVI มีค่าตั้งแต่ $-1$ ถึง $+1$ โดยค่าที่ใกล้ $+1$ บ่งบอกว่าพื้นที่มีพืชพรรณหนาแน่น
\end{ex}

\subsection{การดำเนินการระดับพื้นที่ใกล้เคียง (Neighborhood Operations)}
\index{Neighborhood Operations}

การดำเนินการระดับพื้นที่ใกล้เคียง (Neighborhood Operations) คือ การดำเนินการที่คำนวณค่าของพิกเซลผลลัพธ์จากค่าของพิกเซลเป้าหมายและพิกเซลรอบข้างที่อยู่ภายในหน้าต่างการคำนวณ (Kernel หรือ Window) การดำเนินการระดับพื้นที่ใกล้เคียงมีความสำคัญอย่างยิ่งในการวิเคราะห์ลักษณะภูมิประเทศ การตรวจจับขอบของวัตถุ และการปรับแต่งข้อมูลเชิงพื้นที่

หน้าต่างการคำนวณอาจมีรูปร่างต่างกัน เช่น สี่เหลี่ยมจัตุรัส วงกลม หรือรูปร่างที่กำหนดเอง ขนาดของหน้าต่างกำหนดระยะทางสูงสุดจากพิกเซลเป้าหมายที่จะรวมในการคำนวณ การเลือกขนาดหน้าต่างขึ้นอยู่กับลักษณะของปรากฏการณ์ที่ต้องการวิเคราะห์และความละเอียดเชิงพื้นที่ของข้อมูล

\begin{ex}
การคำนวณค่าเฉลี่ยความสูงในพื้นที่ (Local Mean Elevation) สำหรับพิกเซลที่ตำแหน่ง $(i, j)$ ค่าเฉลี่ยจะคำนวณจากค่าความสูงของพิกเซลเป้าหมายและพิกเซลรอบข้างทั้งหมดในหน้าต่างขนาด $3 \times 3$ หากหน้าต่างมีพิกเซล 9 ตัว ค่าเฉลี่ยจะเป็น

\[
\bar{H}_{ij} = \frac{1}{9} \sum_{m=-1}^{1} \sum_{n=-1}^{1} H_{i+m, j+n}
\]

การคำนวณนี้ทำให้ได้กริดใหม่ที่แสดงค่าเฉลี่ยความสูงในพื้นที่โดยรอบแต่ละพิกเซล ซึ่งช่วยลดผลกระทบของค่าผิดปกติ (Outliers) และสร้างแผนที่ความสูงที่เรียบขึ้น
\end{ex}

\begin{ex}
การคำนวณความลาดชัน (Slope) ความลาดชันเป็นการวัดอัตราการเปลี่ยนแปลงของค่าระดับความสูงในทิศทางที่ลาดที่สุด สามารถคำนวณได้จากค่าความสูงของพิกเซลและพิกเซลรอบข้างในหน้าต่าง $3 \times 3$ โดยใช้สูตรคำนวณความลาดชันดังนี้

\[
S = \arctan \sqrt{\left( \frac{dz}{dx} \right)^2 + \left( \frac{dz}{dy} \right)^2}
\]

โดยที่ $\frac{dz}{dx}$ และ $\frac{dz}{dy}$ คืออัตราการเปลี่ยนแปลงของความสูงในทิศทาง $x$ และ $y$ ตามลำดับ ค่าความลาดชันจะแสดงในหน่วยองศาหรือเปอร์เซ็นต์
\end{ex}

\begin{ex}
การคำนวณค่าเบี่ยงเบนมาตรฐานของความสูง (Local Standard Deviation) ค่าเบี่ยงเบนมาตรฐานเป็นการวัดความไม่สม่ำเสมอของพื้นที่ ค่าที่สูงบ่งบอกว่าพื้นที่มีความแปรปรวนสูง เช่น พื้นที่ภูเขาที่มีหุบเขาสลับกับยอดเขา ค่าเบี่ยงเบนมาตรฐานคำนวณจาก

\[
\sigma_{ij} = \sqrt{ \frac{1}{N-1} \sum_{k=1}^{N} (H_k - \bar{H}_{ij})^2 }
\]

โดยที่ $N$ คือจำนวนพิกเซลในหน้าต่าง และ $\bar{H}_{ij}$ คือค่าเฉลี่ยของความสูงในหน้าต่าง
\end{ex}

\subsection{การดำเนินการระดับโซน (Zonal Operations)}
\index{Zonal Operations}

การดำเนินการระดับโซน (Zonal Operations) คือ การดำเนินการที่คำนวณค่าสถิติสำหรับกลุ่มของพิกเซลที่อยู่ในโซนเดียวกัน โดยโซนถูกกำหนดจากกริดอ้างอิง (Zone Grid) ที่แต่ละพิกเซลมีค่าเป็นรหัสของโซนที่พิกเซลนั้นสังกัด การดำเนินการระดับโซนมักใช้ในการคำนวณสถิติของคุณลักษณะภายในเขตการปกครอง พื้นที่อนุรักษ์ หรือประเภทการใช้ที่ดิน

ผลลัพธ์ของการดำเนินการระดับโซนจะเป็นกริดใหม่ที่แต่ละพิกเซลมีค่าเป็นค่าสถิติของโซนที่พิกเซลนั้นสังกัด หรืออาจเป็นตารางสรุปที่แสดงค่าสถิติสำหรับแต่ละโซนก็ได้ การเลือกรูปแบบผลลัพธ์ขึ้นอยู่กับวัตถุประสงค์ของการวิเคราะห์

\begin{ex}
การคำนวณพื้นที่เฉลี่ยของแต่ละเขตการปกครอง (Mean Area per Zone) พิจารณากริดความสูงและกริดเขตการปกครอง การคำนวณความสูงเฉลี่ยของแต่ละเขตการปกครองจะได้ผลลัพธ์ดังนี้ สำหรับโซนที่มีรหัส $k$ ความสูงเฉลี่ยจะเป็น

\[
\bar{H}_k = \frac{1}{N_k} \sum_{i \in Z_k} H_i
\]

โดยที่ $Z_k$ คือเซตของพิกเซลที่อยู่ในโซน $k$ และ $N_k$ คือจำนวนพิกเซลในโซน $k$
\end{ex}

\begin{ex}
การคำนวณความยาวของแม่น้ำในแต่ละลุ่มน้ำ (Total Stream Length per Watershed) การดำเนินการระดับโซนสามารถใช้ในการรวมค่าของพิกเซลที่เป็นแม่น้ำภายในแต่ละลุ่มน้ำ ผลลัพธ์จะเป็นตารางที่แสดงความยาวรวมของแม่น้ำในแต่ละลุ่มน้ำ ซึ่งมีประโยชน์ในการประเมินทรัพยากรน้ำและการวางแผนการจัดการลุ่มน้ำ
\end{ex}

\begin{ex}
การนับจำนวนพิกเซลในแต่ละโซน (Pixel Count per Zone) การดำเนินการระดับโซนสามารถใช้นับจำนวนพิกเซลในแต่ละโซนเพื่อคำนวณพื้นที่รวมของแต่ละโซน หากทราบขนาดของพิกเซล (เช่น $30 \times 30$ เมตร) สามารถคำนวณพื้นที่ได้โดย

\[
A_k = N_k \times (\Delta x \times \Delta y)
\]

โดยที่ $A_k$ คือพื้นที่ของโซน $k$ และ $\Delta x$, $\Delta y$ คือขนาดของพิกเซลในทิศทาง $x$ และ $y$
\end{ex}

\section{การวิเคราะห์พื้นผิว (Surface Analysis)}
\index{การวิเคราะห์พื้นผิว}
\index{Surface Analysis}

\subsection{แนะนำการวิเคราะห์พื้นผิว}
การวิเคราะห์พื้นผิว (Surface Analysis) เป็นการศึกษาคุณลักษณะและพฤติกรรมของปรากฏการณ์ที่มีการเปลี่ยนแปลงอย่างต่อเนื่องในพื้นที่ ข้อมูลพื้นผิวมักได้จากข้อมูลระดับความสูง (Digital Elevation Model: DEM) ซึ่งเป็นการแสดงค่าความสูงของพื้นผิวโลกในรูปแบบราสเตอร์ การวิเคราะห์พื้นผิวมีการประยุกต์ใช้หลากหลาย เช่น การวางแผนเส้นทาง การประเมินภัยพิบัติ และการจัดการทรัพยากรธรรมชาติ

พื้นผิวทางกายภาพที่พบบ่อยที่สุดคือ ความสูงของภูมิประเทศ ซึ่งสามารถวิเคราะห์คุณลักษณะต่างๆ ได้แก่ ความลาดชัน (Slope) ทิศทางของความลาด (Aspect) เงา (Hillshade) และพื้นที่ที่มองเห็นได้ (Viewshed) นอกจากนี้ยังสามารถวิเคราะห์พื้นผิวอื่นๆ เช่น อุณหภูมิ ความชื้น และความดันบรรยากาศ

การวิเคราะห์พื้นผิวต้องอาศัยข้อมูลระดับความสูงที่มีความละเอียดเพียงพอเพื่อให้ได้ผลลัพธ์ที่ถูกต้อง ความละเอียดของข้อมูลระดับความสูงควรเลือกให้เหมาะสมกับขนาดของลักษณะภูมิประเทศที่ต้องการวิเคราะห์ หากความละเอียดต่ำเกินไป อาจไม่สามารถตรวจจับลักษณะภูมิประเทศขนาดเล็กได้

\subsection{ความลาดชัน (Slope)}
\index{ความลาดชัน}
\index{Slope}

ความลาดชัน (Slope) คือ อัตราการเปลี่ยนแปลงของความสูงต่อหนึ่งหน่วยระยะทางในทิศทางที่ลาดที่สุด ความลาดชันเป็นคุณลักษณะพื้นฐานที่สำคัญของภูมิประเทศ มีการใช้งานหลากหลาย เช่น การประเมินความเสี่ยงดินถล่ม การวางแผนการก่อสร้าง และการจัดหาข้อมูลสำหรับการเกษตร

ความลาดชันสามารถแสดงได้สองรูปแบบ ได้แก่ ความลาดชันในหน่วยองศา (Degree Slope) ที่มีค่าตั้งแต่ $0^\circ$ ถึง $90^\circ$ และความลาดชันในหน่วยเปอร์เซ็นต์ (Percent Slope) ที่มีค่าเป็น $100 \times \tan(\theta)$ โดยที่ $\theta$ คือความลาดชันในหน่วยองศา ความลาดชันในหน่วยเปอร์เซ็นต์จะมีค่าเป็น $0\%$ สำหรับพื้นราบ และเพิ่มขึ้นโดยไม่มีขอบเขตสำหรับพื้นที่ดิ่ง

การคำนวณความลาดชันใช้ค่าความสูงของพิกเซลเป้าหมายและพิกเซลรอบข้างในหน้าต่าง $3 \times 3$ โดยมีสูตรดังนี้

\[
\begin{aligned}
\frac{dz}{dx} &= \frac{(z_{i+1,j+1} + 2z_{i+1,j} + z_{i+1,j-1}) - (z_{i-1,j+1} + 2z_{i-1,j} + z_{i-1,j-1})}{8 \cdot \Delta x} \\
\frac{dz}{dy} &= \frac{(z_{i-1,j-1} + 2z_{i,j-1} + z_{i+1,j-1}) - (z_{i-1,j+1} + 2z_{i,j+1} + z_{i+1,j+1})}{8 \cdot \Delta y}
\end{aligned}
\]

โดยที่ $z_{i,j}$ คือค่าความสูงที่พิกเซล $(i,j)$ และ $\Delta x$, $\Delta y$ คือขนาดของพิกเซล

\begin{ex}
พิจารณาหน้าต่าง $3 \times 3$ ของข้อมูลระดับความสูงที่มีค่าเป็นเมตร ดังแสดงในตารางที่~\ref{tab:slope-calc}

\begin{table}[ht]
\centering
\caption{ตัวอย่างข้อมูลความสูงสำหรับคำนวณความลาดชัน (หน่วย: เมตร)}
\label{tab:slope-calc}
\begin{tabular}{|c|c|c|c|}
\hline
 & \textbf{คอลัมน์ 1} & \textbf{คอลัมน์ 2} & \textbf{คอลัมน์ 3} \\ \hline
\textbf{แถว 1} & 100 & 110 & 105 \\ \hline
\textbf{แถว 2} & 102 & 108 & 115 \\ \hline
\textbf{แถว 3} & 98  & 105 & 112 \\ \hline
\end{tabular}
\end{table}

จากตารางที่~\ref{tab:slope-calc} พิกเซลเป้าหมายอยู่ที่แถว 2 คอลัมน์ 2 มีค่า $z_{2,2} = 108$ เมตร สมมติขนาดพิกเซล $\Delta x = \Delta y = 30$ เมตร คำนวณได้ $\frac{dz}{dx} = \frac{(105 + 2(115) + 112) - (100 + 2(102) + 98)}{8 \times 30} = \frac{447 - 402}{240} = 0.1875$ เมตรต่อเมตร และ $\frac{dz}{dy} = \frac{(100 + 2(102) + 98) - (105 + 2(115) + 112)}{8 \times 30} = \frac{402 - 447}{240} = -0.1875$ เมตรต่อเมตร ความลาดชันจะเป็น $\arctan(\sqrt{0.1875^2 + (-0.1875)^2}) \approx 15.5^\circ$
\end{ex}

\subsection{ทิศทางของความลาด (Aspect)}
\index{ทิศทางของความลาด}
\index{Aspect}

ทิศทางของความลาด (Aspect) คือ ทิศทางที่พื้นที่หันไปหรือทิศทางที่ลาดลงมากที่สุด Aspect มีค่าเป็นมุมที่วัดจากทิศเหนือในทิศทางตามเข็มนาฬิกา มีค่าตั้งแต่ $0^\circ$ (ทิศเหนือ) ถึง $360^\circ$ (กลับมาทิศเหนือ) การทราบทิศทางของความลาดมีความสำคัญในการวิเคราะห์การรับแสงอาทิตย์ การกระจายตัวของพืชพรรณ และการประเมินผลกระทบของลม

ทิศทางของความลาดคำนวณจากอัตราการเปลี่ยนแปลงของความสูงในทิศทาง $x$ และ $y$ โดยใช้สูตร

\[
\phi = 180^\circ + \arctan 2\left( \frac{dz}{dy}, -\frac{dz}{dx} \right)
\]

โดยที่ $\phi$ คือทิศทางของความลาดในหน่วยองศา และ $\arctan 2$ คือฟังก์ชันอาร์แทนเจนต์ที่คืนค่ามุมในช่วง $-180^\circ$ ถึง $180^\circ$ ค่าทิศทางจะถูกปรับให้อยู่ในช่วง $0^\circ$ ถึง $360^\circ$

\begin{ex}
จากตัวอย่างการคำนวณความลาดชันในหัวข้อก่อนหน้า $\frac{dz}{dx} = 0.1875$ และ $\frac{dz}{dy} = -0.1875$ คำนวณทิศทางของความลาดได้ $\phi = 180^\circ + \arctan 2(-0.1875, -0.1875) = 180^\circ + 135^\circ = 315^\circ$ ซึ่งหมายความว่าความลาดหันไปทางทิศตะวันตกเฉียงเหนือ (Northwest)
\end{ex}

\begin{ex}
การประยุกต์ใช้ในการวิเคราะห์การรับแสงอาทิตย์ ทิศทางของความลาดมีผลต่อปริมาณแสงอาทิตย์ที่พื้นที่ได้รับ ในซีกโลกเหนือ �พื้นที่ที่หันไปทางทิศใต้ (Aspect ประมาณ $180^\circ$) จะได้รับแสงอาทิตย์มากกว่าพื้นที่ที่หันไปทางทิศเหนือ ส่งผลให้มีอุณหภูมิสูงกว่าและความชื้นต่ำกว่า ข้อมูลนี้มีประโยชน์ในการวางแผนการเพาะปลูกและการจัดวางอาคาร
\end{ex}

\subsection{เงา (Hillshade)}
\index{เงา}
\index{Hillshade}

เงา (Hillshade) คือ แผนที่ที่แสดงความสว่างของพื้นผิวตามทิศทางของแสงอาทิตย์ การคำนวณเงาใช้ทิศทางและมุมของรังสีแสงอาทิตย์เป็นอินพุต ผลลัพธ์จะเป็นกริดที่แต่ละพิกเซลมีค่าความสว่างตั้งแต่ $0$ ถึง $255$ (สำหรับข้อมูล 8-bit) หรือ $0$ ถึง $1$ (สำหรับข้อมูลเชิงปริมาณปกติ)

การคำนวณเงาใช้สูตร

\[
H = 255 \times \cos(\theta_s) \times \cos(V_{ij}) + 255 \times \sin(\theta_s) \times \sin(V_{ij}) \times \cos(\phi_s - \phi_{ij})
\]

โดยที่ $\theta_s$ คือมุมทิฐิของดวงอาทิตย์ (Solar Zenith Angle) $\phi_s$ คือทิศทางอะซิมัทของดวงอาทิตย์ (Solar Azimuth) และ $V_{ij}$ และ $\phi_{ij}$ คือมุมทิฐิและทิศทางอะซิมัทของพิกเซล $(i,j)$

\begin{ex}
การสร้างแผนที่เงาสำหรับการวิเคราะห์ภูมิประเทศ แผนที่เงาช่วยให้เห็นลักษณะภูมิประเทศได้ชัดเจนยิ่งขึ้น โดยเฉพาะในพื้นที่ที่มีความลาดชันมากและมีหุบเขาลึก การใช้แผนที่เงาร่วมกับข้อมูลอื่น เช่น การใช้ที่ดิน ช่วยให้การตีความข้อมูลเชิงพื้นที่มีประสิทธิภาพมากขึ้น
\end{ex}

\begin{ex}
การประยุกต์ใช้ในการวิเคราะห์ทัศนียภาพ (Visual Impact Analysis) การคำนวณเงาสามารถใช้ในการประเมินว่าอาคารหรือโครงสร้างใหม่จะสร้างเงาตกกระทบพื้นที่ใดบ้างในช่วงเวลาต่างๆ ของวันและฤดูกาล ข้อมูลนี้มีประโยชน์ในการวางแผนผังเมืองและการออกแบบอาคาร
\end{ex}

\subsection{พื้นที่ที่มองเห็นได้ (Viewshed)}
\index{พื้นที่ที่มองเห็นได้}
\index{Viewshed}

พื้นที่ที่มองเห็นได้ (Viewshed) คือ พื้นที่บนพื้นผิวที่สามารถมองเห็นจากตำแหน่งที่กำหนดได้ การวิเคราะห์ Viewshed มีการใช้งานหลากหลาย เช่น การวางตำแหน่งอาคารสูง การวางแผนเส้นทางถนน การจัดตำแหน่งสถานีฐานสัญญาณ และการประเมินทัศนียภาพของภูมิประเทศ

หลักการคำนวณ Viewshed คือการลากเส้นตรงจากตำแหน่งผู้สังเกตไปยังแต่ละพิกเซลในกริด และตรวจสอบว่ามีพื้นที่สูงกว่าเส้นนี้หรือไม่ หากมี แสดงว่าพิกเซลนั้นไม่สามารถมองเห็นได้จากตำแหน่งผู้สังเกต

\begin{ex}
การวิเคราะห์ Viewshed สำหรับการจัดตำแหน่งหอสังเกตการณ์ สมมติต้องเลือกตำแหน่งหอสังเกตการณ์ที่สามารถมองเห็นพื้นที่ได้มากที่สุด การวิเคราะห์ Viewshed จะคำนวณพื้นที่ที่มองเห็นได้จากแต่ละตำแหน่งที่เป็นไปได้ และเลือกตำแหน่งที่มีพื้นที่มองเห็นมากที่สุด
\end{ex}

\begin{ex}
การวิเคราะห์ Viewshed สำหรับการจัดตำแหน่งสถานีฐานสัญญาณโทรศัพท์มือถือ การวิเคราะห์ Viewshed สามารถประเมินพื้นที่ที่จะได้รับสัญญาณจากสถานีฐาน โดยคำนึงถึงอุปสรรคจากภูมิประเทศ ผลลัพธ์ช่วยในการวางแผนจำนวนและตำแหน่งของสถานีฐานให้ครอบคลุมพื้นที่บริการ
\end{ex}

\section{การแปลงข้อมูล Vector-Raster}
\index{การแปลง Vector-Raster}

\subsection{ความสำคัญของการแปลงข้อมูล}
ข้อมูลเชิงพื้นที่มีอยู่สองรูปแบบหลักคือ Vector และ Raster การวิเคราะห์บางประเภทต้องใช้ข้อมูลรูปแบบใดรูปแบบหนึ่งโดยเฉพาะ ดังนั้นการแปลงข้อมูลระหว่างสองรูปแบบนี้จึงมีความจำเป็นอย่างยิ่ง การแปลงข้อมูลไม่ได้เป็นเพียงการเปลี่ยนรูปแบบการจัดเก็บข้อมูล แต่ยังเกี่ยวข้องกับการตัดสินใจเชิงวิเคราะห์ว่าข้อมูลใดเหมาะสมกับการวิเคราะห์ใด

ข้อมูล Vector เหมาะสำหรับการแสดงวัตถุที่มีขอบเขตชัดเจน เช่น อาคาร ถนน และเขตการปกครอง ข้อมูล Raster เหมาะสำหรับการแสดงปรากฏการณ์ที่มีการเปลี่ยนแปลงอย่างต่อเนื่อง เช่น อุณหภูมิ ความสูง และการกระจายตัวของสารเคมี

\subsection{การแปลง Raster เป็น Vector (Raster to Vector Conversion)}
\index{Raster to Vector}

การแปลง Raster เป็น Vector (R2V) คือ การแปลงข้อมูลราสเตอร์ให้เป็นข้อมูลเวกเตอร์ที่ประกอบด้วยจุด เส้น และรูปหลายเหลี่ยม การแปลงนี้มักใช้ในการสร้างแผนที่เวกเตอร์จากภาพถ่ายดาวเทียมหรือภาพถ่ายทางอากาศ

กระบวนการแปลง Raster เป็น Vector ประกอบด้วยหลายขั้นตอน ได้แก่ การปรับปรุงคุณภาพของภาพ (Image Enhancement) การแบ่งส่วนภาพ (Image Segmentation) การจัดหมวดหมู่ (Classification) และการสร้างเวกเตอร์ (Vectorization) การปรับปรุงคุณภาพของภาพช่วยเพิ่มความชัดเจนของวัตถุ การแบ่งส่วนภาพแบ่งภาพออกเป็นพื้นที่ที่มีลักษณะคล้ายกัน การจัดหมวดหมู่กำหนดประเภทของวัตถุในแต่ละพื้นที่ และการสร้างเวกเตอร์สร้างเส้นขอบเขตของวัตถุ

\begin{ex}
การแปลงข้อมูลการใช้ที่ดินเป็นเวกเตอร์ ข้อมูลการใช้ที่ดินในรูปแบบราสเตอร์ที่ได้จากการจัดหมวดหมู่ภาพถ่ายดาวเทียมสามารถแปลงเป็นรูปหลายเหลี่ยมที่แทนเขตการใช้ที่ดินแต่ละประเภทได้ ผลลัพธ์เป็นข้อมูลเวกเตอร์ที่สามารถใช้ในการวิเคราะห์เชิงพื้นที่ร่วมกับข้อมูลเวกเตอร์อื่น เช่น เส้นถนนและแหล่งน้ำ
\end{ex}

\begin{ex}
การแปลงเส้นชั้นความสูงเป็นเวกเตอร์ เส้นชั้นความสูง (Contour Lines) ที่ได้จากข้อมูลระดับความสูงสามารถสร้างเป็นข้อมูลเวกเตอร์ที่ประกอบด้วยเส้นที่เชื่อมต่อกันแทนตำแหน่งที่มีความสูงเท่ากัน ข้อมูลเส้นชั้นความสูงเวกเตอร์มีประโยชน์ในการทำแผนที่ภูมิประเทศและการวิเคราะห์การมองเห็น
\end{ex}

\subsection{การแปลง Vector เป็น Raster (Vector to Raster Conversion)}
\index{Vector to Raster}

การแปลง Vector เป็น Raster (V2R) คือ การแปลงข้อมูลเวกเตอร์ให้เป็นข้อมูลราสเตอร์ การแปลงนี้มักใช้ในการเตรียมข้อมูลสำหรับการวิเคราะห์ราสเตอร์หรือการทำแผนที่แบบ Raster

วิธีการแปลง Vector เป็น Raster หลักมีหลายวิธี ได้แก่ วิธีจุดศูนย์กลาง (Centroid Method) ที่กำหนดค่าของพิกเซลจากค่าของวัตถุที่จุดศูนย์กลางของพิกเซล วิธีเส้นข้าม (Line Intersection Method) ที่พิกเซลใดมีเส้นของวัตถุข้ามผ่านจะถูกกำหนดให้เป็นค่าของวัตถุนั้น และวิธีการกระจายค่า (Area Weighting Method) ที่กระจายค่าของวัตถุไปยังพิกเซลตามสัดส่วนพื้นที่ที่วัตถุครอบคลุม

\begin{ex}
การแปลงข้อมูลเขตการปกครองเป็นราสเตอร์ ข้อมูลเขตการปกครองในรูปแบบรูปหลายเหลี่ยมสามารถแปลงเป็นกริดที่แต่ละพิกเซลมีค่าเป็นรหัสของเขตการปกครองที่พิกเซลนั้นตกอยู่ กริดผลลัพธ์สามารถใช้ในการวิเคราะห์ระดับโซน (Zonal Operations) ร่วมกับข้อมูลราสเตอร์อื่น เช่น การคำนวณพื้นที่เฉลี่ยของแต่ละเขตการปกครอง
\end{ex}

\begin{ex}
การแปลงข้อมูลจุดเป็นราสเตอร์ด้วยวิธีการประมาณค่า (Interpolation) ข้อมูลจุด เช่น สถานีวัดอุณหภูมิ สามารถแปลงเป็นกริดที่แสดงการกระจายของอุณหภูมิโดยใช้วิธีการประมาณค่า เช่น Inverse Distance Weighting (IDW) หรือ Kriging ผลลัพธ์จะเป็นข้อมูลราสเตอร์ที่แสดงค่าอุณหภูมิสำหรับทุกตำแหน่งในพื้นที่
\end{ex}

\section{การประยุกต์ใช้ในงานด้านทรัพยากรธรรมชาติและสิ่งแวดล้อม}
\index{การประยุกต์ใช้ทรัพยากรธรรมชาติ}

การวิเคราะห์ข้อมูลราสเตอร์มีการประยุกต์ใช้อย่างกว้างขวางในงานด้านทรัพยากรธรรมชาติและสิ่งแวดล้อม การประยุกต์ใช้เหล่านี้มีความสำคัญในการจัดการทรัพยากรอย่างยั่งยืน การประเมินผลกระทบสิ่งแวดล้อม และการวางแผนการใช้ที่ดิน

การประยุกต์ใช้หลักในงานด้านทรัพยากรธรรมชาติและสิ่งแวดล้อมประกอบด้วย การประเมินความเหมาะสมของที่ดิน การวิเคราะห์ความเสี่ยงดินถล่ม การประเมินศักยภาพของที่ดินสำหรับการเพาะปลูก การวิเคราะห์การกระจายตัวของสัตว์ป่า และการติดตามการเปลี่ยนแปลงของสิ่งแวดล้อม

การใช้ Map Algebra ร่วมกับข้อมูลหลายชั้น (Multi-criteria Analysis) ช่วยให้สามารถประเมินความเหมาะสมของพื้นที่ตามเกณฑ์หลายประการพร้อมกัน การวิเคราะห์นี้มักใช้ในการจัดหาพื้นที่สำหรับวัตถุประสงค์เฉพาะ เช่น การเลือกพื้นที่สำหรับโรงไฟฟ้าพลังงานแสงอาทิตย์หรือการจัดหาพื้นที่เพาะปลูกใหม่

\subsection{การประเมินความเหมาะสมของที่ดิน (Land Suitability Analysis)}
\index{การประเมินความเหมาะสมของที่ดิน}

การประเมินความเหมาะสมของที่ดินเป็นกระบวนการที่ใช้การวิเคราะห์ราสเตอร์หลายชั้นในการประเมินว่าพื้นที่ใดเหมาะสมสำหรับวัตถุประสงค์ใด กระบวนการนี้ประกอบด้วยการกำหนดเกณฑ์ การให้คะแนนแต่ละเกณฑ์ การรวมคะแนน และการจัดระดับความเหมาะสม

เกณฑ์ที่ใช้ในการประเมินความเหมาะสมของที่ดินขึ้นอยู่กับวัตถุประสงค์ ตัวอย่างเช่น การประเมินความเหมาะสมสำหรับการเพาะปลูกข้าวอาจพิจารณาเกณฑ์ ได้แก่ ความลาดชัน (ควรน้อยกว่า $8^\circ$) ความสูง (ควรต่ำกว่า 500 เมตร) ประเภทดิน และความใกล้ของแหล่งน้ำ

\begin{ex}
การประเมินความเหมาะสมของที่ดินสำหรับการเพาะปลูกมันสำปะหลัง พิจารณาเกณฑ์ดังนี้ ความลาดชันไม่เกิน $15^\circ$ ปริมาณน้ำฝนรายปีระหว่าง $1{,}000$ ถึง $2{,}000$ มิลลิเมตร และดินเป็นดินร่วนหรือดินทรายร่วน การใช้ Map Algebra สามารถสร้างแผนที่ความเหมาะสมได้โดย

\[
S = (S_{slope} \geq 0.8) \land (S_{rain} \geq 0.8) \land (S_{soil} \geq 0.8)
\]

โดยที่ $S_{slope}$, $S_{rain}$, และ $S_{soil}$ คือค่าความเหมาะสมของแต่ละเกณฑ์ที่ปรับให้อยู่ในช่วง $0$ ถึง $1$
\end{ex}

\begin{ex}
การประเมินความเหมาะสมของที่ดินสำหรับโครงการอนุรักษ์ป่า การประเมินความเหมาะสมต้องพิจารณาทั้งความเหมาะสมทางนิเวศวิทยาและความเป็นไปได้ทางเศรษฐกิจและสังคม เกณฑ์ทางนิเวศวิทยาอาจรวมถึงความหลากหลายทางชีวภาพ ความสมบูรณ์ของระบบนิเวศ และความเชื่อมต่อกับพื้นที่อนุรักษ์อื่น
\end{ex}

\subsection{การวิเคราะห์ความเสี่ยงดินถล่ม (Landslide Risk Analysis)}
\index{การวิเคราะห์ความเสี่ยงดินถล่ม}

การวิเคราะห์ความเสี่ยงดินถล่มใช้การวิเคราะห์ราสเตอร์หลายชั้นในการระบุพื้นที่ที่มีความเสี่ยงสูงต่อการเกิดดินถล่ม ปัจจัยที่มีผลต่อการเกิดดินถล่มประกอบด้วย ความลาดชัน ประเภทหิน ปริมาณน้ำฝน การระบายน้ำ และการใช้ที่ดิน

การสร้างแผนที่ความเสี่ยงดินถล่มใช้การรวมปัจจัยต่างๆ โดยกำหนดน้ำหนักของแต่ละปัจจัยตามความสำคัญ การวิเคราะห์นี้มีประโยชน์ในการวางแผนการใช้ที่ดิน การเตือนภัยล่วงหน้า และการจัดการภัยพิบัติ

\begin{ex}
การคำนวณดัชนีความเสี่ยงดินถล่ม (Landslide Susceptibility Index: LSI) ใช้สูตร

\[
\text{LSI} = \sum_{i=1}^{n} (w_i \times f_i)
$$

โดยที่ $w_i$ คือน้ำหนักของปัจจัย $i$ และ $f_i$ คือค่าปกติของปัจจัย $i$ ตัวอย่างเช่น น้ำหนักของความลาดชันอาจเป็น $0.4$ และน้ำหนักของปริมาณน้ำฝนอาจเป็น $0.3$ ผลลัพธ์จะเป็นกริดที่มีค่าความเสี่ยงตั้งแต่ $0$ ถึง $1$ โดยค่าที่สูงกว่าหมายถึงความเสี่ยงที่สูงกว่า
\end{ex}

\section{บทสรุป}
\index{สรุป}

การวิเคราะห์ข้อมูลราสเตอร์เป็นเครื่องมือสำคัญในระบบสารสนเทศภูมิศาสตร์ที่มีการประยุกต์ใช้อย่างกว้างขวางในงานด้านทรัพยากรธรรมชาติและสิ่งแวดล้อม ในบทนี้ได้กล่าวถึงหลักการพื้นฐานของโครงสร้างข้อมูลราสเตอร์ ซึ่งเป็นรูปแบบการจัดเก็บข้อมูลเชิงพื้นที่ที่ใช้กริดของพิกเซลในการแทนค่าความเข้มข้นหรือคุณลักษณะ

พีชคณิตแผนที่ (Map Algebra) เป็นชุดของการดำเนินการทางคณิตศาสตร์และตรรกะที่ใช้ในการวิเคราะห์ข้อมูลราสเตอร์ การดำเนินการแบ่งได้เป็น 4 ประเภทตามขอบเขตการคำนวณ ได้แก่ Local Operations ที่ทำงานกับพิกเซลเดี่ยว Neighborhood Operations ที่ทำงานกับพิกเซลรอบข้าง Zonal Operations ที่ทำงานกับกลุ่มพิกเซลในโซนเดียวกัน และ Global Operations ที่ทำงานกับพิกเซลทั้งหมด

การวิเคราะห์พื้นผิว (Surface Analysis) ใช้ข้อมูลระดับความสูงในการคำนวณคุณลักษณะต่างๆ ของภูมิประเทศ ได้แก่ ความลาดชัน (Slope) ทิศทางของความลาด (Aspect) เงา (Hillshade) และพื้นที่ที่มองเห็นได้ (Viewshed) คุณลักษณะเหล่านี้มีการประยุกต์ใช้หลากหลายในงานด้านวิศวกรรม การวางแผน และการจัดการทรัพยากร

การแปลงข้อมูลระหว่างรูปแบบ Vector และ Raster มีความจำเป็นอย่างยิ่งในการวิเคราะห์ข้อมูลเชิงพื้นที่ เนื่องจากการวิเคราะห์บางประเภทต้องใช้ข้อมูลในรูปแบบใดรูปแบบหนึ่งโดยเฉพาะ การเลือกรูปแบบข้อมูลที่เหมาะสมขึ้นอยู่กับลักษณะของปรากฏการณ์ที่ต้องการวิเคราะห์และวิธีการวิเคราะห์ที่จะใช้

ความรู้ความเข้าใจในการวิเคราะห์ข้อมูลราสเตอร์เป็นพื้นฐานสำคัญสำหรับการประยุกต์ใช้ในงานด้านทรัพยากรธรรมชาติและสิ่งแวดล้อม เช่น การประเมินความเหมาะสมของที่ดินและการวิเคราะห์ความเสี่ยงดินถล่ม ซึ่งมีความสำคัญต่อการจัดการทรัพยากรอย่างยั่งยืนและการวางแผนการใช้ที่ดิน

\section{แบบฝึกหัดท้ายบท}
\index{แบบฝึกหัด}

\begin{enumerate}
    \item อธิบายความแตกต่างระหว่างข้อมูลราสเตอร์และข้อมูลเวกเตอร์ พร้อมยกตัวอย่างประเภทข้อมูลที่เหมาะสมกับแต่ละรูปแบบ

    \item จากข้อมูลระดับความสูงในรูปแบบราสเตอร์ขนาด $5 \times 5$ จงคำนวณความลาดชันและทิศทางของความลาดของพิกเซลกลาง โดยมีค่าความสูงดังนี้ (หน่วย: เมตร ขนาดพิกเซล: 30 เมตร)

\[
\begin{bmatrix}
100 & 105 & 110 & 108 & 105 \\
102 & 108 & 115 & 112 & 107 \\
98  & 105 & 112 & 118 & 110 \\
95  & 100 & 108 & 115 & 112 \\
90  & 95  & 102 & 110 & 108
\end{bmatrix}
\]

    \item อธิบายความแตกต่างระหว่าง Local Operations, Neighborhood Operations, Zonal Operations และ Global Operations ใน Map Algebra พร้อมยกตัวอย่างการใช้งานของแต่ละประเภท

    \item กำหนดข้อมูลการใช้ที่ดินและข้อมูลความลาดชัน จงอธิบายขั้นตอนการหาพื้นที่ที่เหมาะสมสำหรับการเพาะปลูกพืชไร่ โดยพืชไร่ต้องการความลาดชันไม่เกิน $10^\circ$ และต้องอยู่ในเขตการใช้ที่ดินประเภทเกษตรกรรม

    \item อธิบายกระบวนการคำนวณ Viewshed และยกตัวอย่างการประยุกต์ใช้ในงานด้านวิศวกรรมหรือการวางแผน

    \item เปรียบเทียบวิธีการแปลง Vector เป็น Raster และ Raster เป็น Vector พร้อมระบุข้อดีและข้อจำกัดของแต่ละวิธี

    \item กำหนดข้อมูลระดับความสูง ข้อมูลปริมาณน้ำฝน และข้อมูลประเภทดิน จงออกแบบการวิเคราะห์ความเสี่ยงดินถล่มโดยใช้ Map Algebra พร้อมระบุสูตรและเกณฑ์การให้คะแนน
\end{enumerate}

\section{เอกสารอ้างอิงประจำบท}
\index{เอกสารอ้างอิง}

\begin{enumerate}
    \item Longley, P.A., Goodchild, M.F., Maguire, D.J., \& Rhind, D.W. (2015). Geographic Information Science and Systems. John Wiley \& Sons.
    \item Tomlin, C.D. (1990). Geographic Information Systems and Cartographic Modeling. Prentice Hall.
    \item Burrough, P.A., \& McDonnell, R.A. (1998). Principles of Geographical Information Systems. Oxford University Press.
    \item Chang, K. (2015). Introduction to Geographic Information Systems. McGraw-Hill Education.
    \item DeMers, M.N. (2009). Fundamentals of Geographic Information Systems. John Wiley \& Sons.
\end{enumerate}
